\section{Finite analytic representations and reusable identities}
\label{sec:analytic}

\subsection{Rational enclosures and elementary operations}
For $z=x+\ii y$ use
$\normone z=|x|+|y|$ and $|z|_\infty=\max(|x|,|y|)$.
These rational norms satisfy
\[
 |z|_\infty\le |z|\le\normone z\le2|z|,
 \qquad \normone{zw}\le\normone z\,\normone w.
\]
A complex ball is a rational centre $c$ and a rational error $r\ge0$,
with meaning $\normone{z-c}\le r$. For addition the radii add. For
multiplication an admissible radius is
\begin{equation}
 \normone c\,r_d+\normone d\,r_c+r_cr_d.
 \label{eq:ballproduct}
\end{equation}
Unless an estimate is explicitly in $\normone{\cdot}$, analytic
absolute-value bounds below use the complex Euclidean modulus. An
implementation that inserts one of them into an L1 radius first
multiplies it by two. The logarithm and exponential error estimates
used by the ball evaluator are direct L1 estimates.
If a centre is rounded to a dyadic rational, add its exact rational
rounding distance to the radius. Division is used only with a supplied
positive lower bound, or with an exactly nonzero rational denominator.

For rational $v>0$, reduce $v$ to $2^ku$, $1\le u<2$. The number of
reductions is bounded by the bit lengths of the numerator and denominator.
With $y=(u-1)/(u+1)$,
\begin{equation}
 \log u=2\sum_{j=0}^{L-1}\frac{y^{2j+1}}{2j+1}+r_L,
 \quad
 |r_L|\le\frac{2|y|^{2L+1}}{(2L+1)(1-y^2)}.
 \label{eq:log-series}
\end{equation}
Here $|y|\le1/3$. A linear bound on $L$ in the requested precision,
plus the bit length of $k$, therefore suffices. For a complex rational
$z$ of norm at most an integer $S\ge1$, the exponential series through
degree $d\ge2S$ has error
\begin{equation}
 \normone{\exp z-\sum_{j=0}^d z^j/j!}\le 2^{4S-d}.
 \label{eq:exp-series}
\end{equation}
Indeed the $2S$th term is at most $\exp S\le4^S$, and the subsequent
term ratios are at most $1/2$. Input perturbation of norm $r\le1$
costs at most $r4^{S+1}$. Both estimates are finite geometric-series
inequalities with their numerical parameters fixed.

The identity $\pi=4\arctan1$ with $2^{p+3}$ alternating terms is one
deliberately slow primitive recursive approximation to error $2^{-p}$.
Faster familiar formulas may be substituted after proving their error
contract. These choices make every displayed expression involving
$\pi$, real logarithms, exponentials, sine or cosine a specified
primitive recursive rational approximation.

\subsection{Euler summation and continuation}
Let $\{x\}=x-\lfloor x\rfloor$. Integrating on each interval
$[n,n+1]$ and summing gives, initially for $\sigma>1$,
\begin{equation}
 \zeta(s)=\frac{s}{s-1}
       -s\int_1^\infty\{x\}x^{-s-1}\,dx,
 \qquad \sigma>0,\quad s\ne1.
 \label{eq:euler-basic}
\end{equation}
The integral on the right has tail at most
$X^{-\sigma}/\sigma$. It consequently supplies a continuation with
an explicit convergence modulus on every region
$\sigma\ge a>0$, $|s|\le H$, $|s-1|\ge b>0$.
Differentiating $j$ times is justified by the bound
\begin{equation}
 \int_1^\infty(\log x)^j x^{-a-1}\,dx=\frac{j!}{a^{j+1}}.
 \label{eq:logmoments}
\end{equation}
One can avoid the word differentiation in the arithmetic proof:
expand $\exp(-h\log x)$ to the required finite degree, bound its
remainder, and integrate the resulting finite identities.

For wider disks and for the executable example we need the full finite
Euler--Maclaurin expression. Put
$(s)_j=s(s+1)\cdots(s+j-1)$ and let
$\widetilde B_j(x)=B_j(\{x\})$ denote the periodic Bernoulli polynomial.
For integers $N,m\ge1$ and $\sigma+2m>1$,
\begin{align}
 \zeta(s)={}&\sum_{n=1}^{N-1}n^{-s}
       +\frac{N^{1-s}}{s-1}+\frac12N^{-s}\notag\\
 &+\sum_{j=1}^m\frac{B_{2j}}{(2j)!}(s)_{2j-1}N^{1-s-2j}
       +R_{m,N}(s),\label{eq:em}\\
 R_{m,N}(s)={}&-\frac{(s)_{2m}}{(2m)!}
       \int_N^\infty\widetilde B_{2m}(x)x^{-s-2m}\,dx.
 \label{eq:em-remainder}
\end{align}
To prove this, integrate by parts on each unit interval, using
$B_j'=jB_{j-1}$ and $B_j(1)=B_j(0)$ for $j\ne1$. All inner endpoint
terms cancel. Repeat $2m$ times and then take a finite upper endpoint
to its explicitly bounded tail. This is a finite induction on the
number of integrations and on the number of unit cells.

The Bernoulli numbers are computed by
\[
 B_0=1,\qquad
 B_n=-\frac1{n+1}\sum_{j<n}\binom{n+1}{j}B_j.
\]
Thus
\begin{equation}
 C_m=\frac1{(2m)!}\sum_{j=0}^{2m}\binom{2m}{j}|B_j|
 \label{eq:bernoulli-bound}
\end{equation}
is a rational upper bound for
$\sup_{0\le x\le1}|B_{2m}(x)|/(2m)!$, obtained simply by bounding
each monomial. Consequently
\begin{equation}
 |R_{m,N}(s)|\le
 C_m\prod_{j<2m}(|\sigma+j|+|t|)
 \frac{N^{1-\sigma-2m}}{\sigma+2m-1}.
 \label{eq:em-bound}
\end{equation}
For the rational $\normone{\cdot}$-ball implementation, multiply the
right side by two. The rotation in $x^{-it}$ need not preserve the
$\normone{\cdot}$ norm, although it preserves the complex modulus.
For $\sigma\ge0$, replace the final factor by
$N^{1-2m}/(2m-1)$. This wholly rational bound is the one implemented in
\code{code/rational_zeta.py}. It requires no numerical library value of
zeta and no assumed Riemann hypothesis.

On a bounded rational box with $\sigma\ge-A$, choose $m=A+2$.
Then the exponent in the tail is at most $-A-3$. A power of two larger
than the rational coefficient bound times $2^{p+10}$ is already a
sufficient $N$ for error $2^{-p-10}$. Approximate the finite elementary
terms with a further budget for their rational coefficient magnitudes.
This gives a fixed primitive recursive approximation on every compact
box apart from $s=1$. To approximate the entire numerator everywhere,
multiply the finite formula symbolically by $s-1$ and cancel the
explicit pole term before evaluation. In particular, no evaluation of
an undefined zeta value at $s=1$ is performed.

\subsection{The functional-equation dependency}
For completeness we spell out the precise identities and quantitative
content reused from the included proof~\cite{RepoFE}. Define
\begin{align}
 \psi(x)&=\sum_{n\ge1}\exp(-\pi n^2x),\notag\\
 J(s)&=\int_1^\infty\psi(x)x^{s/2-1}\,dx,\notag\\
 \xi(s)&=\frac12+\frac{s(s-1)}2\bigl(J(s)+J(1-s)\bigr).
 \label{eq:theta-xi}
\end{align}
The theta relation
$1+2\psi(x)=x^{-1/2}(1+2\psi(1/x))$ follows there from a finite
Fej\'er calculation for the Gaussian, with Gaussian and Fourier tails
bounded before a limit is taken. Splitting the Mellin integral at $1$
gives~\eqref{eq:xi-definition} on $\sigma>1$ and its continuation.
Equation~\eqref{eq:theta-xi} gives both symmetries directly, with no
choice of a branch of the logarithm.

Here is a useful part of its explicit approximation schedule. On
$|\Re s|,|\Im s|\le B$, $B\ge2$, write
$\lambda(n)=\min\{j\le n:n\le2^j\}$ for $n\ge1$. Given $p$, put
\begin{align*}
 X&=p+\lambda(B!)+4,& K&=p+(B+1)\lambda(X)+4,\\
 D&=2(B+2)X^B,& q&=p+\lambda(XD)+4,&
 r&=p+\lambda(X)+4.
\end{align*}
Truncate the $x$ integral at $X$, the theta sum at $K$, use a left
mesh of width $2^{-q}$, and evaluate its finite summands so that the
error of their entire sum is at most $2^{-r}$ at each sample.
The four error bounds are
\begin{equation}
 2B!2^{-2X},\quad
 4X^{B+1}2^{-3(K+1)^2},\quad
 XD2^{-q},\quad X2^{-r}.
 \label{eq:theta-errors}
\end{equation}
Their sum is less than $2^{-p}$. These bounds follow respectively
from $\psi(x)\le2e^{-3x}$, a geometric estimate in the theta index,
a derivative bound for the integrand, and the sum of sample errors.
For example, $\psi\le1$ and
$|\psi'|\le4\sum_{n\ge1}n^2 8^{-n}=288/343<1$ justify the
displayed derivative budget. To obtain error $2^{-p}$ for $\xi$ itself,
first enclose both $s$ and $1-s$ in one common integer box of radius $B$,
and apply the integral schedule at precision
$p+\lambda(4B^2)+3$. This absorbs the two integral errors and the
multiplication by $s(s-1)/2$.
The same approximation parameters used for $s$ and $1-s$ make the
finite xi expression exactly symmetric.

The reciprocal gamma function is represented by
\begin{equation}
 G_N(z)=zN^{-z}\prod_{n=1}^N(1+z/n),\qquad G(z)=1/\Gamma(z).
 \label{eq:gamma-product}
\end{equation}
For integer $D\ge1$, $|z|\le D$, the included proof supplies
\begin{align}
 |G_N(z)|,|G(z)|&\le M_D:=D4^{D+D^2},\notag\\
 |G_N(z)-G(z)|&\le C_D/N,
 \quad C_D=2M_D(D+D^2),\quad N\ge2(D+D^2).
 \label{eq:gamma-modulus}
\end{align}
Away from the nonpositive integers it also gives a positive product
lower bound. The shift identity is $G(z)=zG(z+1)$.
These statements, their theta/Mellin identification, and the Gaussian
normalization $\Gamma(1/2)=\sqrt\pi$ are the reusable premises.
They are written proofs with finite numerical contracts, not completed
Lean theorems. Their source archive is preserved unchanged.

\subsection{From analytic formulas to rational polynomials}
Every use of a polynomial approximation below has a numerical error
requested in advance. On a closed disk strictly inside a domain of
holomorphy, a Taylor truncation is one construction: a bound $M_*$ on
a concentric radius $R_*$ gives
\[
 \sup_{|z-c|\le r}\left|f(z)-\sum_{j=0}^d a_j(z-c)^j\right|
 \le\frac{M_* (r/R_*)^{d+1}}{1-r/R_*}.
\]
The coefficients are obtained by finite Cauchy quadrature with an
explicit derivative bound. Alternatively one can use the fixed finite
Euler--Maclaurin or theta expressions above and replace each exponential
by its Taylor polynomial after approximating its real constants.
This second route avoids importing a general Taylor theorem: its degree,
coefficients and error budget follow by finite sum and product estimates.

All required disk sizes, larger-box bounds, degrees and working
precisions are fixed primitive recursive expressions in the input data.
The choice may be extremely inefficient. Efficiency is not used to
infer provability; the finite error estimates and induction invariants
are the reason this representation is suitable for $\ipra$.
